\DocumentMetadata{
  pdfversion=2.0,
  pdfstandard=ua-2,
  lang=en-US,
  tagging=on,
  tagging-setup={math/setup=mathml-SE,tabsorder=structure}
}
\documentclass[11pt,letterpaper]{article}
\usepackage[margin=0.68in,includefoot]{geometry}
\usepackage{newcomputermodern}
\usepackage{amsmath,amscd,mathtools}
\providecommand{\square}{\mdlgwhtsquare}
\usepackage[hidelinks]{hyperref}
\hypersetup{
  pdftitle={Combinatorics PhD Exam, August 2005},
  pdfauthor={University of Florida Department of Mathematics},
  pdfsubject={Graduate mathematics examination},
  pdfdisplaydoctitle=true,
  bookmarksopen=true
}
\setcounter{secnumdepth}{0}
\setlength{\parindent}{0pt}
\setlength{\parskip}{0.28em}
\setlength{\emergencystretch}{3em}
\setlength{\leftmargini}{2.15em}
\setlength{\leftmarginii}{2.65em}
\setlength{\leftmarginiii}{3em}
\renewcommand{\labelenumi}{\textbf{\theenumi.}}
\pagestyle{empty}
\raggedbottom
\allowdisplaybreaks
\newenvironment{examproblems}[1]{%
  \begin{enumerate}%
  \setcounter{enumi}{#1}%
  \setlength{\topsep}{0.35em}%
  \setlength{\partopsep}{0pt}%
  \setlength{\itemsep}{0.48em}%
  \setlength{\parsep}{0.18em}%
}{\end{enumerate}}
\newenvironment{examsubparts}{%
  \begin{enumerate}%
  \setlength{\topsep}{0.18em}%
  \setlength{\partopsep}{0pt}%
  \setlength{\itemsep}{0.16em}%
  \setlength{\parsep}{0.1em}%
}{\end{enumerate}}
\newenvironment{examinstructions}{%
  \par\begingroup\itshape
}{%
  \par\endgroup\vspace{0.18em}
}
\makeatletter
\renewcommand\section{\@startsection{section}{1}{0pt}%
  {-1.25ex plus -.25ex minus -.1ex}%
  {0.55ex plus .1ex}%
  {\normalfont\large\bfseries}}
\renewcommand{\@maketitle}{%
  \newpage\null\vskip -1.2em%
  {\footnotesize\bfseries
    \href{https://www.ufl.edu/}{University of Florida}, \href{https://math.ufl.edu/}{Department of Mathematics}\par}%
  \vskip 0.18em%
  \tagstructbegin{tag=Title}%
    {\LARGE\bfseries\href{https://gma.math.ufl.edu/exams/combinatorics-phd/}{Combinatorics PhD Exam}, August 2005\par}%
  \tagstructend%
  \vskip 0.35em%
  \tagmcbegin{artifact}%
    \hrule height 1pt%
  \tagmcend%
  \vskip 0.55em%
}
\makeatother
\title{Combinatorics PhD Exam, August 2005}
\author{}
\date{}
\begin{document}
\maketitle
\thispagestyle{empty}
\begin{examinstructions}
Show your work.
\end{examinstructions}
\begin{examproblems}{0}
\item
This problem has two parts.
\begin{examsubparts}
\item[\textnormal{a.}]
There are~\(r\) black and \(n-r\) white balls in an urn. They are removed one at a time without replacement. What is the probability that exactly~\(k\) drawings are required to get a white ball?
\item[\textnormal{b.}]
Use your result to conclude that 
\[
\sum_{k=1}^{r}\frac{(r)_k}{(n-1)_k}=\frac{r}{n-r}.
\]
\end{examsubparts}
\item
Each of~\(n\) people is to be mailed an envelope containing a letter and a bill. How many ways \(Q_n\) are there of placing the~\(n\) letters and~\(n\) bills into~\(n\) addressed envelopes so that no envelope contains both the correct letter and bill?
\item
Let \(P_n\) be the total number of~\(k\)-permutations of~\(n\) for various~\(k\), that is, 
\[
P_n=\sum_{k=0}^{n}(n)_k,\qquad n=0,1,\ldots
\]
 Show that 
\[
P(t)=\sum_{n=0}^{\infty}P_n\frac{t^n}{n!}=(1-t)^{-1}e^t
\]
 and use this to show that 
\[
P_n=nP_{n-1}+1,\qquad n=1,2,\ldots,\quad P_0=1.
\]
\item
Define the binary Hamming code \(H(r)\) of length \(2^r-1\). Show that \(H(r)\) is an exactly single error correcting code and that \(H(r)\) is a perfect code. Determine the number of codewords of weight 3 in \(H(r)\).
\item
Show that the number of partitions of a number~\(n\) into exactly~\(m\) parts is equal to the number of partitions of \(n-m\) into no more than~\(m\) parts.
\item
An order on the set of ordered pairs of non-negative integers is defined by \((a_1,a_2)\leq(b_1,b_2)\) if \(a_i\leq b_i\) for \(i=1,2\). Find the Möbius function of this poset.
\item
Determine for which values of~\(m\) and~\(n\) the complete bipartite graph \(K_{mn}\) is. Answer the same three questions for the~\(n\)-cube. (Recall that the n-cube \(Q_n\) is defined as the graph whose vertices are the set of all binary sequences of length~\(n\), where two vertices are adjacent if the corresponding sequences differ by exactly one digit.)
\begin{examsubparts}
\item[\textnormal{(a)}]
planar
\item[\textnormal{(b)}]
Eulerian
\item[\textnormal{(c)}]
Hamiltonian.
\end{examsubparts}
\item
Let~\(G\) be a triangle-free graph with~\(n\) vertices, minimum degree~\(k\) and girth~\(g\). Prove that \(g\leq 2n/k\).

\par
Hint: For an appropriate cycle~\(C\) count the number of edges from~\(C\) to \(G\mathbin{\backslash}C\) in two ways.
\end{examproblems}
\end{document}
